% ============================================================
%  Báo cáo cải tiến: CCNet-4Class Training Pipeline
%  Mô tả các thay đổi trong scripts/new_feature/ nhằm bám sát
%  Wang et al. (IEEE TIV 2023) và cải thiện val mIoU từ ~0.79
%  lên mục tiêu 85–90%.
% ============================================================

\documentclass[12pt,a4paper]{article}

\usepackage[utf8]{inputenc}
\usepackage[T5]{fontenc}
\usepackage[vietnamese,english]{babel}
\usepackage{geometry}
\geometry{margin=2.5cm, top=3cm, bottom=3cm}
\usepackage{amsmath,amssymb}
\usepackage{graphicx}
\usepackage{booktabs}
\usepackage{array}
\usepackage{multirow}
\usepackage{xcolor}
\usepackage{hyperref}
\usepackage{caption}
\usepackage{float}
\usepackage{enumitem}
\usepackage{setspace}
\usepackage{fancyhdr}
\usepackage{listings}
\usepackage{tikz}
\usetikzlibrary{shapes.geometric, arrows.meta, positioning, calc}

\definecolor{codebg}{RGB}{248,248,248}
\definecolor{codeframe}{RGB}{200,200,200}
\definecolor{keyword}{RGB}{0,0,180}
\definecolor{comment}{RGB}{100,100,100}
\definecolor{string}{RGB}{160,0,0}

\lstset{
  backgroundcolor=\color{codebg},
  frame=single,
  rulecolor=\color{codeframe},
  basicstyle=\ttfamily\small,
  keywordstyle=\color{keyword}\bfseries,
  commentstyle=\color{comment}\itshape,
  stringstyle=\color{string},
  numbers=left,
  numberstyle=\tiny\color{gray},
  breaklines=true,
  language=Python,
  showstringspaces=false,
  tabsize=4,
}

\pagestyle{fancy}
\fancyhf{}
\rhead{CCNet-4Class: Cải tiến Training Pipeline}
\lhead{SAC Paper Reproduction}
\rfoot{\thepage}

\hypersetup{
  colorlinks=true,
  linkcolor=blue,
  urlcolor=blue,
}

\onehalfspacing

% ─────────────────────────────────────────────────────────────
\begin{document}

\begin{titlepage}
  \centering
  \vspace*{2cm}
  {\LARGE \textbf{Cải Tiến Pipeline Huấn Luyện CCNet-4Class}}\\[0.5cm]
  {\large Tái hiện Wang et al. ``Semantic-Aware Video Compression\\
  for Automotive Cameras'', IEEE TIV 2023}\\[1.5cm]

  \begin{tabular}{ll}
    \textbf{Phạm vi thay đổi:} & \texttt{scripts/new\_feature/} \\
    \textbf{Các file cập nhật:} & \texttt{ccnet\_4class.py}, \texttt{train.py}, \texttt{dataset.py} \\
    \textbf{Vấn đề ban đầu:} & val mIoU $\approx 0.79$ (mục tiêu paper: $\geq 0.87$) \\
    \textbf{Mục tiêu sau cải tiến:} & val mIoU $85\%$–$90\%$ \\
  \end{tabular}

  \vfill
  \today
\end{titlepage}

\tableofcontents
\newpage

% ============================================================
\section{Tổng quan vấn đề}
% ============================================================

Trong quá trình tái hiện bài báo SAC (Wang et al., IEEE TIV 2023), model CCNet-4Class
được huấn luyện trên tập Cityscapes chỉ đạt val mIoU $\approx 0.79$, thấp hơn đáng kể
so với mức kỳ vọng $\geq 0.87$ mà paper ghi nhận (Table IV: SA-X265 đạt mIoU 87.48\%
trên frame đã nén). Phân tích nguyên nhân xác định ba vấn đề kỹ thuật trong code:

\begin{enumerate}[leftmargin=*, label=\arabic*.]
  \item \textbf{Không có LR scheduler} — Adam với learning rate cố định $3\times10^{-4}$
        suốt 47 epoch. Mô hình plateau sớm khoảng epoch 25–30.
  \item \textbf{Feature map quá thô (stride-32)} — Backbone ResNet101 giữ nguyên stride,
        khiến feature map sau \texttt{layer4} chỉ còn $16\times32$ pixels (tại input
        $512\times1024$). Sau đó upsample $\times32$ lên output — gây mất chi tiết không gian
        nghiêm trọng.
  \item \textbf{Augmentation yếu} — Chỉ dùng horizontal flip, thiếu scale augmentation
        và color jitter.
\end{enumerate}

Ước tính ảnh hưởng lên mIoU:

\begin{table}[H]
  \centering
  \caption{Phân tích nguyên nhân mIoU thấp}
  \begin{tabular}{lcc}
    \toprule
    \textbf{Vấn đề} & \textbf{Ảnh hưởng ước tính} & \textbf{File} \\
    \midrule
    Không có LR scheduler & $-3$ đến $-5\%$ mIoU & \texttt{train.py} \\
    Stride-32 (không dilated) & $-2$ đến $-4\%$ mIoU & \texttt{ccnet\_4class.py} \\
    Augmentation yếu & $-1$ đến $-2\%$ mIoU & \texttt{dataset.py} \\
    \midrule
    \textbf{Tổng} & \boldmath$-6$ đến $-11\%$ mIoU & — \\
    \bottomrule
  \end{tabular}
\end{table}

% ============================================================
\section{Thay đổi 1 — Kiến trúc Dilated Backbone (\texttt{ccnet\_4class.py})}
% ============================================================

\subsection{Vấn đề: Stride-32 features}

ResNet101 tiêu chuẩn có chuỗi stride:

\[
  \text{Input} \xrightarrow{\times4} \text{layer1}
               \xrightarrow{\times2} \text{layer2}
               \xrightarrow{\times2} \text{layer3}
               \xrightarrow{\times2} \text{layer4}
  \quad\Rightarrow\quad \text{stride tổng} = 32
\]

Với input $512\times1024$, feature map sau \texttt{layer4} chỉ còn $16\times32$.
Sau khi upsample $\times32$ về kích thước gốc, thông tin không gian bị mất rất nhiều.
Đây là lý do chính các mô hình segmentation như DeepLab, CCNet gốc (ICCV 2019) và
PSPNet đều sử dụng \textbf{dilated (atrous) convolution} trên \texttt{layer3} và
\texttt{layer4}: giữ nguyên stride $\times8$, tăng receptive field bằng dilation.

\subsection{Giải pháp: Dilated layer3 và layer4}

Thêm helper \texttt{\_make\_dilated} để chuyển đổi mà \textbf{không thay đổi shape
của weights} (checkpoint cũ vẫn tương thích):

\begin{lstlisting}[caption={Hàm chuyển đổi sang dilated convolution}]
@staticmethod
def _make_dilated(layer: nn.Sequential, dilation: int) -> None:
    for block in layer:
        # 3x3 conv trong Bottleneck là conv2
        block.conv2.dilation = (dilation, dilation)
        block.conv2.padding  = (dilation, dilation)
        block.conv2.stride   = (1, 1)
        # Xoá stride ở projection shortcut (block đầu tiên)
        if block.downsample is not None:
            block.downsample[0].stride = (1, 1)
\end{lstlisting}

Áp dụng trong \texttt{\_\_init\_\_}:

\begin{lstlisting}[caption={Kích hoạt dilated trong constructor}]
self.layer3 = bb.layer3   # dilation=2 -> stays H/8 x W/8
self.layer4 = bb.layer4   # dilation=4 -> stays H/8 x W/8

self._make_dilated(self.layer3, dilation=2)
self._make_dilated(self.layer4, dilation=4)
\end{lstlisting}

\subsection{Kết quả}

\begin{table}[H]
  \centering
  \caption{So sánh feature map trước và sau khi dilated}
  \begin{tabular}{lccc}
    \toprule
    \textbf{Version} & \textbf{Feature map sau layer4} & \textbf{Upsample} & \textbf{Spatial info} \\
    \midrule
    Cũ (stride-32) & $16\times32$ & $\times32$ & Rất thô \\
    Mới (stride-8)  & $64\times128$ & $\times8$  & Tốt hơn $4\times$ \\
    \bottomrule
  \end{tabular}
\end{table}

\subsection{CrissCross Attention — Viết lại cho stride-8}

Với feature map $64\times128 = 8192$ vị trí, attention đầy đủ $N\times N$ yêu cầu
$8192^2 \approx 67M$ phần tử/batch — không khả thi trên GPU thực tế. Giải pháp:
\textbf{row-wise + column-wise attention} tách biệt, giảm xuống
$O(N\cdot(H+W))$ thay vì $O(N^2)$:

\begin{lstlisting}[caption={Row + column attention thay thế full N×N}]
# Row-wise: mỗi pixel attend với toàn bộ pixel cùng hàng
Q_r = Q.permute(0,2,1,3).reshape(B*H, c8, W)
K_r = K.permute(0,2,1,3).reshape(B*H, c8, W)
V_r = V.permute(0,2,1,3).reshape(B*H,  C, W)
attn_r = F.softmax(torch.bmm(Q_r.permute(0,2,1), K_r)*scale, dim=-1)
out_r  = torch.bmm(V_r, attn_r.permute(0,2,1))
out_r  = out_r.reshape(B,H,C,W).permute(0,2,1,3)

# Column-wise: mỗi pixel attend với toàn bộ pixel cùng cột
Q_c = Q.permute(0,3,1,2).reshape(B*W, c8, H)
...
out = self.gamma * (out_r + out_c) + x
\end{lstlisting}

Hai lần áp dụng CCA stacked (cc1, cc2) vẫn xấp xỉ full attention như trong paper
gốc CCNet (ICCV 2019).

\begin{table}[H]
  \centering
  \caption{Chi phí bộ nhớ attention tại feature map $64\times128$}
  \begin{tabular}{lcc}
    \toprule
    \textbf{Loại attention} & \textbf{Phần tử (batch=4)} & \textbf{Khả thi?} \\
    \midrule
    Full $N\times N$ & $4 \times 8192^2 \approx 268M$ & Không \\
    Row-wise ($H$ ma trận $W\times W$) & $4\times64\times128^2 \approx 4M$ & Có \\
    Col-wise ($W$ ma trận $H\times H$) & $4\times128\times64^2 \approx 2M$ & Có \\
    \bottomrule
  \end{tabular}
\end{table}

% ============================================================
\section{Thay đổi 2 — Poly LR Scheduler \& Separate LR (\texttt{train.py})}
% ============================================================

\subsection{Vấn đề: LR cố định}

Với Adam lr $= 3\times10^{-4}$ không đổi suốt 47 epoch, mô hình thường plateau
ở epoch 25–30 và không cải thiện thêm. Hầu hết các mô hình segmentation hiện đại
(CCNet, DeepLab, PIDNet) dùng \textbf{poly LR decay}.

\subsection{Giải pháp 1: Poly LR Decay}

\[
  \text{lr}(e) = \text{lr}_\text{base} \times \left(1 - \frac{e}{E_\text{max}}\right)^{0.9}
\]

trong đó $e$ là epoch hiện tại, $E_\text{max}$ là tổng số epoch. Power $= 0.9$ là
giá trị tiêu chuẩn trong các paper segmentation.

\begin{lstlisting}[caption={Poly LR decay}]
def poly_lr(base_lr: float, epoch: int, max_epochs: int,
            power: float = 0.9) -> float:
    return base_lr * (1.0 - epoch / max_epochs) ** power

# Trong training loop — cập nhật mỗi epoch:
for pg in optim.param_groups:
    pg["lr"] = poly_lr(pg["initial_lr"], epoch, args.epochs)
\end{lstlisting}

\begin{figure}[H]
  \centering
  \begin{tikzpicture}
    \begin{scope}[xscale=5.5, yscale=3]
      % Axes
      \draw[->] (0,0) -- (1.08,0) node[right] {Epoch / $E_{\max}$};
      \draw[->] (0,0) -- (0,1.15) node[above] {lr / lr$_{\text{base}}$};
      % Poly curve
      \draw[blue, thick, domain=0:1, samples=60]
        plot (\x, {(1-\x)^0.9});
      % Flat baseline
      \draw[red, dashed, thick] (0,1) -- (1,1);
      % Labels
      \node[blue] at (0.6, 0.55) {Poly ($p=0.9$)};
      \node[red]  at (0.5, 1.08) {Cũ: cố định};
      % Ticks
      \foreach \x in {0, 0.5, 1}
        \draw (\x, 0.02) -- (\x, -0.02) node[below] {\x};
      \foreach \y in {0, 0.5, 1}
        \draw (0.01, \y) -- (-0.01, \y) node[left] {\y};
    \end{scope}
  \end{tikzpicture}
  \caption{So sánh LR schedule cũ (cố định) và mới (poly decay)}
\end{figure}

\subsection{Giải pháp 2: Separate LR — Backbone vs Head}

Backbone đã được pretrain trên ImageNet, chỉ cần fine-tune nhẹ.
Head (reduction + CCA + seg\_head) được khởi tạo ngẫu nhiên, cần LR cao hơn.

\begin{lstlisting}[caption={Separate learning rate cho backbone và head}]
backbone_params = (
    list(model.stem.parameters())
    + list(model.layer1.parameters())
    + list(model.layer2.parameters())
    + list(model.layer3.parameters())
    + list(model.layer4.parameters())
)
head_params = (
    list(model.reduction.parameters())
    + list(model.cc1.parameters())
    + list(model.cc2.parameters())
    + list(model.seg_head.parameters())
)

optim = torch.optim.Adam([
    {"params": backbone_params,
     "lr": BASE_LR * 0.1,        # backbone: 10x thấp hơn
     "initial_lr": BASE_LR * 0.1},
    {"params": head_params,
     "lr": BASE_LR,              # head: full LR
     "initial_lr": BASE_LR},
], betas=(0.9, 0.999), weight_decay=1e-5)
\end{lstlisting}

\begin{table}[H]
  \centering
  \caption{Learning rate theo epoch (base\_lr = $3\times10^{-4}$, 60 epochs)}
  \begin{tabular}{lccc}
    \toprule
    \textbf{Thành phần} & \textbf{Epoch 0} & \textbf{Epoch 30} & \textbf{Epoch 59} \\
    \midrule
    Backbone & $3.0\times10^{-5}$ & $1.6\times10^{-5}$ & $7.5\times10^{-7}$ \\
    Head     & $3.0\times10^{-4}$ & $1.6\times10^{-4}$ & $7.5\times10^{-6}$ \\
    \bottomrule
  \end{tabular}
\end{table}

\subsection{Tăng epochs: 47 $\rightarrow$ 60}

Với poly LR decay, các epoch đầu mang LR cao (hội tụ nhanh) và các epoch cuối
mang LR rất thấp (tinh chỉnh). 47 epoch thường không đủ để tận dụng phase tinh chỉnh này.
60 epochs là mức khuyến nghị phổ biến trong các paper segmentation Cityscapes.

\subsection{Cải thiện resume checkpoint}

Code cũ: \texttt{strict=True} khi load state dict. Nếu kiến trúc thay đổi
(dilated vs non-dilated), việc load sẽ fail hoàn toàn. Code mới dùng \texttt{strict=False}
và báo warning về missing/unexpected keys — cho phép partial load (warm-start).

% ============================================================
\section{Thay đổi 3 — Augmentation (\texttt{dataset.py})}
% ============================================================

\subsection{Vấn đề: Chỉ dùng horizontal flip}

Augmentation duy nhất là \texttt{FLIP\_LEFT\_RIGHT} với $p=0.5$. Thiếu:
\begin{itemize}
  \item Scale augmentation — mô hình không học được invariance với kích thước đối tượng
  \item Color jitter — dễ overfit màu sắc cụ thể của Cityscapes
\end{itemize}

\subsection{Giải pháp: 4-step augmentation pipeline}

Thứ tự áp dụng tương tự DeepLab / CCNet:

\begin{lstlisting}[caption={Pipeline augmentation mới (training only)}]
def _augment(self, image, label, H, W):
    # 1. Random scale: [0.75, 1.5]x base_size, clip min = base_size
    scale = random.uniform(0.75, 1.5)
    new_h = max(int(H * scale), H)
    new_w = max(int(W * scale), W)
    image = image.resize((new_w, new_h), Image.BILINEAR)
    label = label.resize((new_w, new_h), Image.NEAREST)

    # 2. Random crop: cat về đúng (H, W)
    top  = random.randint(0, new_h - H)
    left = random.randint(0, new_w - W)
    image = image.crop((left, top, left + W, top + H))
    label = label.crop((left, top, left + W, top + H))

    # 3. Random horizontal flip (p=0.5)
    if random.random() > 0.5:
        image = image.transpose(Image.FLIP_LEFT_RIGHT)
        label = label.transpose(Image.FLIP_LEFT_RIGHT)

    # 4. Color jitter (chỉ áp dụng cho image, không áp dụng label)
    image = self._color_jitter(image)
    return image, label
\end{lstlisting}

Tham số color jitter: \texttt{brightness=0.4, contrast=0.4, saturation=0.4, hue=0.1} —
đây là giá trị phổ biến trong các paper segmentation Cityscapes.

\begin{table}[H]
  \centering
  \caption{So sánh augmentation cũ và mới}
  \begin{tabular}{lcc}
    \toprule
    \textbf{Augmentation} & \textbf{Cũ} & \textbf{Mới} \\
    \midrule
    Resize cố định $512\times1024$ & \checkmark & \checkmark \\
    Horizontal flip ($p=0.5$) & \checkmark & \checkmark \\
    Random scale {[}0.75, 1.5{]} & $\times$ & \checkmark \\
    Random crop về $512\times1024$ & $\times$ & \checkmark \\
    Color jitter & $\times$ & \checkmark \\
    \bottomrule
  \end{tabular}
\end{table}

% ============================================================
\section{Hướng dẫn chạy lại}
% ============================================================

\subsection{Lệnh train lại từ đầu (khuyến nghị)}

\begin{lstlisting}[language=bash, caption={Chạy trên GPU thuê}]
cd /path/to/sac_project/scripts

# Train lại từ đầu với kiến trúc và schedule mới
python new_feature/train.py --reset --epochs 60 --batch-size 4

# Theo dõi kết quả
tail -f ../outputs/new_feature/training/train_metrics.csv
\end{lstlisting}

\textbf{Lưu ý:} Bắt buộc dùng \texttt{--reset} vì checkpoint cũ được train với
stride-32. Load checkpoint cũ vào model stride-8 sẽ cho feature distribution sai —
cần train lại từ pretrained ImageNet backbone.

\subsection{Kỳ vọng kết quả}

\begin{table}[H]
  \centering
  \caption{Kỳ vọng val mIoU theo từng cải tiến}
  \begin{tabular}{lcc}
    \toprule
    \textbf{Cấu hình} & \textbf{val mIoU (ước tính)} & \textbf{Ghi chú} \\
    \midrule
    Baseline (trước cải tiến) & $\approx 0.79$ & stride-32, lr cố định \\
    + Poly LR & $\approx 0.82$–$0.84$ & — \\
    + Dilated stride-8 & $\approx 0.85$–$0.87$ & Cải tiến lớn nhất \\
    + Augmentation & $\approx 0.87$–$0.90$ & Mục tiêu paper \\
    \midrule
    \textbf{Paper reference} & $\mathbf{\geq 0.87}$ & Table IV, SA-X265 \\
    \bottomrule
  \end{tabular}
\end{table}

\subsection{Đánh giá sau khi train xong}

\begin{lstlisting}[language=bash, caption={Chạy evaluate trên toàn bộ val set}]
python new_feature/evaluate.py --num-frames 500 --preset fast
\end{lstlisting}

Kết quả lưu tại \texttt{outputs/new\_feature/evaluation/<timestamp>/}:
\begin{itemize}
  \item \texttt{table2\_cityscapes.csv} — PSNR, SSIM, SA-PSNR, SA-SSIM
  \item \texttt{table4\_segmentation.csv} — mIoU, iIoU trên frame đã nén
  \item \texttt{results\_summary.txt} — bảng so sánh với paper reference
\end{itemize}

% ============================================================
\section{Tóm tắt thay đổi}
% ============================================================

\begin{table}[H]
  \centering
  \caption{Danh sách đầy đủ các thay đổi}
  \renewcommand{\arraystretch}{1.3}
  \begin{tabular}{p{3.5cm} p{4cm} p{5.5cm}}
    \toprule
    \textbf{File} & \textbf{Thay đổi} & \textbf{Lý do} \\
    \midrule
    \texttt{ccnet\_4class.py} &
      Dilated layer3 (×2), layer4 (×4) &
      Feature map: $16\times32 \to 64\times128$; upsample $\times32 \to \times8$ \\
    \cmidrule{2-3}
    & CrissCross Attention viết lại thành row+col &
      Full $N\times N$ ($N=8192$) không khả thi với stride-8; row+col tương đương \\
    \midrule
    \texttt{train.py} &
      Poly LR decay (power=0.9) &
      LR cố định plateau sớm; poly LR là standard trong dense-prediction \\
    \cmidrule{2-3}
    & Separate LR: backbone $\times0.1$ vs head $\times1$ &
      Backbone pretrained — fine-tune chậm; head random init — train nhanh \\
    \cmidrule{2-3}
    & Epochs: $47 \to 60$ &
      Cần thêm epoch để tận dụng phase LR thấp \\
    \cmidrule{2-3}
    & \texttt{strict=False} khi load checkpoint &
      Cho phép warm-start khi kiến trúc thay đổi \\
    \midrule
    \texttt{dataset.py} &
      Random scale {[}0.75, 1.5{]} + random crop &
      Scale invariance; standard trong segmentation training \\
    \cmidrule{2-3}
    & Color jitter &
      Tránh overfit màu sắc; tăng robustness \\
    \bottomrule
  \end{tabular}
\end{table}

\end{document}
